\subsection*{A consequence for integral ideals}
In every number field $K$, the set of nonzero integral ideals whose
prime-ideal exponents belong to $G$ admits a strict improvement in
natural density by norm, preserving avoidance even for fractional-ideal
ratios. This gives a qualitative answer to the natural-density aspect
of Best, Huan, McNew, Miller, Powell, Tor and Weinstein's Question 5.1
for quadratic fields
(\href{https://arxiv.org/abs/1412.0999}{their paper, section 5}).
Their question asks how much nongreedy sets can improve existing bounds.
This argument gives a positive improvement on the greedy lower bound
for natural density. Sharp field-specific improvements and the extremal
upper density are separate questions.

Indeed, fix a prime ideal $\mathfrak l$ of norm $Q$ and choose a finite
set $B$ of other prime ideals. In the preceding construction replace
$2$ by $\mathfrak l$, $q$ by $\mathfrak q$, and $P$ by
$\mathfrak l\prod_{\mathfrak q\in B}\mathfrak q$. Every combinatorial
step is unchanged by unique factorization of ideals. Give an ideal
reciprocal mass $1/N\mathfrak a$. In units of $1/NP$, the additions
have weight $Q^{-1}\sum_{\mathfrak q\in B}1/N\mathfrak q$, while the
deletions have weight at most
\[
  \frac{Q}{1-Q^{-1}}
  \left(1+\sum_{\mathfrak q\in B}
    \frac1{N\mathfrak q\,(N\mathfrak q-1)}\right).
\]
The sum in parentheses is bounded uniformly in $B$ since
$\sum_{\mathfrak q}(N\mathfrak q)^{-2}<\infty$.
In contrast, $\sum_{\mathfrak q}1/N\mathfrak q$ diverges:
the logarithm of the Euler product for $\zeta_K(s)$ tends to infinity
as real $s\downarrow1$, and terms from prime-ideal powers of exponent
at least two have bounded total. Thus the gain is positive for some
finite $B$.

To justify natural norm-density, the count $J_K(x)$ of integral ideals
of norm at most $x$ satisfies $J_K(x)\sim c_Kx$ with $c_K>0$
(\href{https://doi.org/10.1016/j.exmath.2006.08.001}{Murty and Van Order}).
Consequently multiples of an ideal $\mathfrak d$ have density
$1/N\mathfrak d$, and $J_K(x)\ll_K x$ for all $x\ge1$.
Inclusion--exclusion yields the local factors for finitely many prime
ideals. The $\mathfrak q^2$-multiple tail has upper density bounded
by a constant times $\sum_{N\mathfrak q>Y}(N\mathfrak q)^{-2}$.
The global density, its positivity and the strict improvement now
follow by the same convergent smooth-part sum.

The density here is the limit of the number of selected ideals of norm at most $x$ divided by $J_K(x)$. We count ideals; associate generators in a ring with infinitely many units would require a separate counting convention.
